% Original: PDF strana 10, donji slajd.
\slajdid{02-s020}
\begin{frame}[label=02-s020]{Invertori i opterećenje izvora signala}
% element: 02-s020:e01
Da li su dostupne i prave i komplementne vrednosti signala? \alert{U većini slučajeva nisu.}\\
Raspolažemo samo pravim vrednostima; komplemente dobijamo invertorima.
\par\vspace{0mm}
\begin{columns}[T,onlytextwidth]
\begin{column}{.65\textwidth}\centering
% element: 02-s020:e02
\begin{tikzpicture}[DEoznaka,x=10mm,y=3.2mm]
\foreach \y/\low/\lab in {9.3/8.1/C,6.7/5.5/B,4.1/2.9/A}{
 \node[not gate US,draw,minimum width=7mm,minimum height=5mm] (n\lab) at (1.6,\y) {};
 \draw[DEvod] (0,\y) node[left] {\(\lab\)}--(n\lab.input);
 \draw[DEvod] (n\lab.output)--(6.2,\y);
 \draw[DEvod] (.6,\y)--(.6,\low)--(6.2,\low);
 \node[junction] at (.6,\y) {};
}
\def\DEtriPrviY{1.0}
\def\DEtriDrugiY{-2.8}
\def\DEtriUlazi{1/3/9.3,1/2/5.5,1/1/2.9,2/3/8.1,2/2/6.7,2/1/2.9,3/3/8.1,3/2/5.5,3/1/4.1}
% Položaji se mogu lokalno prilagoditi, uz iste priključke i polaritete.
\ifdefined\DEtriPrviY\else\def\DEtriPrviY{1.7}\fi
\ifdefined\DEtriDrugiY\else\def\DEtriDrugiY{-1.3}\fi
\ifdefined\DEtriUlazi\else
\def\DEtriUlazi{1/3/7.9,1/2/5.2,1/1/3.3,2/3/7.1,2/2/6.0,2/1/3.3,3/3/7.1,3/2/5.2,3/1/4.1}
\fi
\foreach \i/\x in {1/3.1,2/4.4,3/5.7}{
 \node[and gate US,draw,logic gate inputs=nnn,rotate=-90,minimum width=8mm,minimum height=8mm] (g\i) at (\x,\DEtriPrviY) {};
}
\foreach \i/\j/\y in \DEtriUlazi{
 \draw[DEvod] (g\i.input \j)--(g\i.input \j |- 0,\y);
 \node[junction] at (g\i.input \j |- 0,\y) {};
}
\node[or gate US,draw,logic gate inputs=nnn,rotate=-90,minimum width=8mm,minimum height=8mm] (z) at (4.4,\DEtriDrugiY) {};
\coordinate (DEsopPut) at ($(g2.output)!.5!(z.input 2)$);
\draw[DEvod] (g1.output)--(g1.output |- DEsopPut)-|(z.input 3);
\draw[DEvod] (g2.output)--(z.input 2);
\draw[DEvod] (g3.output)--(g3.output |- DEsopPut)-|(z.input 1);
\draw[DEvod] (z.output)--++(0,-.35) node[below] {\(F\)};

% Ista tri opterećena ulaza kao u originalu; obojeni stvarni vodovi ne uvode nove veze.
\draw[DEvod,draw=DEcrvena,line width=1.1pt] (0,4.1)--(nA.input);
\draw[DEvod,draw=DEcrvena,line width=1.1pt] (.6,4.1)--(.6,2.9)--(g2.input 1 |- 0,2.9);
\draw[DEvod,draw=DEcrvena,line width=1.1pt] (g1.input 1 |- 0,2.9)--(g1.input 1);
\draw[DEvod,draw=DEcrvena,line width=1.1pt] (g2.input 1 |- 0,2.9)--(g2.input 1);
\node[junction,fill=DEcrvena] at (.6,4.1) {};
\node[junction,fill=DEcrvena] at (g1.input 1 |- 0,2.9) {};
\node[junction,fill=DEcrvena] at (g2.input 1 |- 0,2.9) {};
\end{tikzpicture}

\end{column}
\begin{column}{.32\textwidth}
% element: 02-s020:e03
\Oznaka{Primer: ulaz \(A\)}
Izvor signala \(A\) opterećen je sa tri ulaza logičkih kola.\par\vspace{2mm}
\alert{Sa stanovišta realnih logičkih kola to nije dobro.}
\par\vspace{3mm}
\(F=\bar CBA+C\bar BA+CB\bar A\)
\end{column}
\end{columns}
\end{frame}
